\chapter{RENCANA PELAKSANAAN}
\label{chap:rencana}
\begin{petunjuk}
Bab Rencana Pelaksanaan yang merupakan bab penutup Laporan Proposal Tugas Akhir digunakan untuk mendeskripsikan rencana pelaksanaan berupa jadwal dan risiko-risiko yang mungkin dihadapi dan rencana mitigasinya. Tujuan bab ini adalah:
\begin{itemize}
  \item Mahasiswa memiliki rencana yang jelas mengenai pelaksanaan tugas akhir.
  \item Mahasiswa mengenali risiko-risiko pelaksanaan tugas akhir yang mungkin dihadapi dan sudah menyiapkan diri untuk mengantisipasi risiko tersebut.
\end{itemize}
\end{petunjuk}

\section{Jadwal}
\begin{petunjuk}
Cantumkan jadwal pengerjaan tugas akhir dan dilengkapi dengan uraiannya.
\end{petunjuk}

\section{Risiko}
\begin{petunjuk}
Cantumkan 5 risiko tertinggi yang mungkin dihadapi dalam pengerjaan tugas akhir. Risiko yang dicantumkan dapat merupakan risiko dari sisi teknis, risiko dari sisi operasional, risiko dari metode yang dipilih, dan sebagainya. Cantumkan pula rencana mitigasi dari risiko-risiko tersebut.
\end{petunjuk}
